\chapter{动画片二：Chain 怎样决定切一刀}
\label{ch:N-chain-cartoon}

\section{Hom 还留着``神秘铃''，Chain 换成自行车锁}

\begin{metaphor}[为什么换锁]
查表铃声有点贵、有点魔法。Chain 改成：每来一字节就拧一下密码锁。\\
拧到对齐（低几位全 0）——才打开第二、第三关。
\end{metaphor}

\section{三道关剧本}

\begin{enumerate}
  \item 拧锁（每个字节都要拧，很便宜）；
  \item 对齐了吗？没对齐 $\to$ 继续拧；对齐了才往后；
  \item 珠串边变了吗？（和 Hom 同款图画）；
  \item （默认）有没有环？有环才剪。
\end{enumerate}

\begin{figure}[htbp]
\centering
\begin{tikzpicture}[font=\small]
  \node[softbox=Gen4Color] (a) {拧锁};
  \node[softbox=WarnAmber, right=0.5cm of a] (b) {对齐?};
  \node[softbox=WarnAmber, right=0.5cm of b] (c) {珠变?};
  \node[softbox=WarnAmber, right=0.5cm of c] (d) {有环?};
  \node[softbox=CutRed, right=0.5cm of d] (e) {剪};
  \draw[-{Stealth}, thick] (a)--(b)--(c)--(d)--(e);
  \draw[-{Stealth}, thick, dashed] (b.south)--++(0,-0.5)-|
    node[pos=0.2,below,font=\tiny]{否：继续拧} (a.south);
\end{tikzpicture}
\caption{时间花在拧锁上；对齐很少见，贵检查才发生。}
\label{fig:N-chain-flow}
\end{figure}

\section{量化：把颜色弄粗糙}

字节可以先右移几位再当颜色，叫量化。\\
粗糙后颜色种类变少，珠串图案更钝 —— 也是一种旋钮。

\section{并行找刀（可选）}

大文件切成路段同时找，但最终只认\textbf{最左边}那一刀。\\
像多人搜宝藏，谁先在更左的位置喊到，右边就停手。

\section{零基础对照}

\begin{center}
\begin{tabular}{@{}ll@{}}
\toprule
术语 & 动画片里 \\
\midrule
LFSR / $S$ & 密码锁当前刻度 \\
strideOK & 锁对齐了 \\
ChainPrimaryDelta & 珠边变了 \\
beta1 & 有环 \\
stride log & 对齐稀密度旋钮 \\
\bottomrule
\end{tabular}
\end{center}
